% !TEX root = ../main.tex
\section{엔도스랭크와 AWP 나란히 보기}
\label{ch:appendix-er-awp}

\ref{ch:results}장은 이 논문의 모든 점수를 함께 보고하고, 각 페이지랭크(PageRank)를 그것이 매끄럽게 다듬는(평활하는) 원시 차수(화살표를 그냥 센 수)와 견주어 읽어요. 이 부록은 한 층만 걷는 두 점수, 곧 엔도스랭크(EndorseRank)와 AWP(적응형 가중 페이지랭크) 사이의 직접 비교를 모아 두었어요. 같은 기간에서의 가족 평균(비슷한 지표 묶음에 대한 평균값)과 짝지은 차이, 대리 지표(대신 재는 값)별 계수, 그리고 기간 밖에서의 짝지은 차이예요. 이 비교들은 하나도 미리 정해 두지 않았고, 어느 것도 이 논문의 주장을 담지 않아요. 두 점수는 서로 다른 화살표 위를 걷고, 서로 다른 구성개념(점수가 재려고 하는 생각)을 위해 만들어졌어요. 그래서 점검마다 앞서는 점수는 거의 언제나, 그 점검이 세는 화살표를 쓰는 점수예요. 두 점수가 다루는 범위도 달라요. 매칭 지갑 가운데 966개에는 허락 화살표가 하나도 닿지 않고, 381개에는 송금 화살표가 하나도 닿지 않아요. 또 홀드아웃(나중 기간을 떼어 두고 맞혀 보는 시험) 스펜더 1,335개 가운데 318개는 동결일(점수를 고정한 날) 전에 송금 화살표가 없어요. 같은 기간에서 엔도스랭크의 순서는 대부분 동점이 정해요(\ref{sec:rank-divergence}절). 두 홀드아웃 기간 모두에서 두 점수는 저마다 자기가 다듬는 원시 차수보다 낮아요(\ref{sec:holdout-results}절). 그래서 둘 사이에 차이가 있다고 해서, 어느 쪽이든 자기 차수에 무언가를 더한다는 것을 보여 주지는 않아요.

\dissertationSubheading[sec:er-awp-same-window]{같은 기간}

표~\ref{tab:alignment-hybrid}는 두 점수의 가족 평균을 구간과 함께 보여 주고, 표~\ref{tab:tie-sensitivity}의 마지막 행은 두 점수가 서로 얼마나 맞는지 보여 줘요. 두 모습은 엇갈려요. AWP는 송금 가족을 가장 가깝게 따라가고($0.612$) 허락 가족에서는 조금 음수예요. 반면 엔도스랭크는 허락 가족을 가장 가깝게 따라가고($0.375$), 송금 가족은 중간 정도로 따라가요($0.330$). 표~\ref{tab:tau-diff}는 같은 지갑 다시 뽑기 표본 위에서 짝지은 차이 다섯 개를 검사해요(\ref{sec:statistical-tests}절). 이 차이들에는 표의 행 순서대로 i부터 v까지 번호를 붙였어요.

% Real BigQuery data -- regenerate with:
%   2-Dissertation-Draft-Works/wallet-reputation-experiments/scripts/run_dissertation_eval.py --real --export-latex
% Generated by export_latex_results.py -- do not edit by hand unless re-export fails
\begin{table}[htbp]
\centering
\caption[Paired bootstrap tests of differences between Kendall $\tau$ coefficients.]{Paired bootstrap tests of differences between Kendall $\tau$ coefficients (matched cohort, $n=5{,}521$). Family entries use the family-mean $\tau$; the same wallet resamples are used for $a$ and $b$, so the interval is on the paired difference. An interval excluding 0 indicates a difference not attributable to sampling variation at the 5\% level. CI = 95\% percentile bootstrap over wallets (400 paired resamples, seed 42).}
\label{tab:tau-diff}
\fitwidth{%
\begin{tabular}{lccccc}
\toprule
Contrast ($a$ minus $b$) & $\tau_a$ & $\tau_b$ & $\Delta\tau$ & 95\% CI of $\Delta\tau$ & $P(\Delta\tau>0)$ \\
\midrule
EndorseRank allowance minus EndorseRank transfer & 0.375 & 0.330 & +0.045 & [-0.001, 0.084] & 0.970 \\
AWP transfer minus EndorseRank transfer & 0.612 & 0.330 & +0.282 & [0.263, 0.299] & 1.000 \\
EndorseRank allowance minus AWP allowance & 0.375 & -0.024 & +0.398 & [0.356, 0.432] & 1.000 \\
AWP sybil-stability minus EndorseRank sybil-stability & 0.278 & 0.231 & +0.047 & [0.031, 0.062] & 1.000 \\
EndorseRank in-approve degree minus EndorseRank in-degree & 0.376 & 0.341 & +0.035 & [-0.015, 0.079] & 0.925 \\
\bottomrule
\end{tabular}
}
\end{table}


다섯 대비(짝지은 차이) 가운데 셋에서 구간이 0을 벗어나요. AWP의 송금 일치도가 엔도스랭크보다 높고($+0.282$), 엔도스랭크의 허락 일치도가 AWP보다 높아요($+0.398$). 그리고 시빌(가짜 주소를 많이 만들어 속이는 공격) 안정성 가족에서는 AWP가 엔도스랭크보다 앞서요($+0.047$). 나머지 둘은 엔도스랭크를 자기 자신과 비교하는데, 둘 다 차이가 보이지 않아요. 엔도스랭크의 허락 일치도가 송금 일치도보다 높다는 것은 보이지 않았어요(대비~i, $+0.045$, 구간 $[-0.001,0.084]$). 받은 허락 수 일치도가 입차수 일치도보다 높다는 것도 보이지 않았어요(대비~v, $+0.035$, $[-0.015,0.079]$). 이 두 ‘차이 없음’은 어느 쪽의 증거로도 읽을 수 없어요. 엔도스랭크의 동점 때문에 $\tau_b$의 천장(얻을 수 있는 가장 큰 값)은 허락 대리 지표에서는 약 $0.38$이고, 동점이 없는 대리 지표에서는 약 $0.57$이에요(\ref{sec:same-window-alignment}절). 그래서 이 계수들은 같은 잣대 위에 있지 않아요. 허락 그래프 밖의 지갑에 0점을 주는 대신 외톨이 점(화살표가 하나도 없는 점)으로 남겨 두면, 대비~i와~v는 $+1.061$과 $+1.090$이 되고 구간도 0보다 훨씬 위에 있어요.

대비~ii와~iii의 부호가 반대인 것은 구성개념에서 나와요. 각 점수는 자기 화살표로 만든 가족에서 앞서요. 그리고 각 가족의 주된 대리 지표에는 그에 맞는 페이지랭크가 다듬는 입차수가 들어 있어요. 허락 가족에서 엔도스랭크가 앞서는 것은 허락을 받는 지갑 131개 덕분이에요. 엔도스랭크는 이 지갑들을 모두 공통 값보다 위로 올리지만, AWP는 이 131개를 자기 순위의 대부분에 걸쳐 흩어 놓아요. 송금 가족에서 AWP가 앞서는 것은 엔도스랭크의 송금 계수와 견준 결과인데, 이 계수는 빠진 지갑을 외톨이 점으로 남겨 두면 음수로 바뀌어요. 두 차이 모두 어느 점수가 평판을 더 잘 재는지는 말해 주지 않아요.

표~\ref{tab:robustness-sample-size}의 단계별 풀 부분 표본은 그 값을 매칭 코호트의 값과 비교할 수 없어요. 이 부분 표본들에서 각 가족 안의 두 점수 순서는 모든 단계에서 같아요. 송금 가족에서는 AWP가 앞서고, 허락 가족에서는 엔도스랭크가 앞서요.

\dissertationSubheading[sec:proxy-family-detail]{대리 지표별 자세한 결과}

가족 평균 하나는 대리 지표별 계수 여러 개를 하나로 줄인 거예요. 표~\ref{tab:alignment-transfer}부터~\ref{tab:alignment-sybil-stability}까지는 엔도스랭크와 AWP에 대해 이 계수들을 대리 지표 하나씩 보여 줘요. 스피어먼의 로($\rho$), 켄달의 타우($\tau$, 두 줄 세우기가 얼마나 같은 순서인지 나타내는 수), 그리고 $\tau$의 95\% 구간이에요. 상관을 구하기 전에, 모든 대리 지표는 값이 클수록 기대되는 평판이 높다는 뜻이 되도록 방향을 맞춰요. 송금 대리 지표는 AWP를 처음 낸 논문이 AWP를 검증할 때 쓴 바로 그 지표들이에요\cend{3}.

\dissertationSubsubheading[sub:transfer-family]{송금 가족}

% Real BigQuery data -- regenerate with:
%   2-Dissertation-Draft-Works/wallet-reputation-experiments/scripts/run_dissertation_eval.py --real --export-latex
% Generated by export_latex_results.py -- do not edit by hand unless re-export fails
\begin{table}[htbp]
\centering
\caption[Domain alignment with transfer proxies.]{Domain alignment with transfer proxies (inbound in-degree and in-value, the validation proxies used for the AWP baseline). Kendall $\tau$ and Spearman $\rho$ between reputation scores and proxy rankings. CI = 95\% percentile bootstrap over wallets (400 paired resamples, seed 42).}
\label{tab:alignment-transfer}
\fitwidth{%
\begin{tabular}{lccc}
\toprule
Proxy & Spearman $\rho$ & Kendall $\tau$ & 95\% CI of $\tau$ \\
\midrule
\multicolumn{4}{l}{\textit{EndorseRank}} \\
In-degree & 0.418 & 0.341 & [0.319, 0.363] \\
In-value (sum inbound) & 0.396 & 0.319 & [0.300, 0.340] \\
\addlinespace
\multicolumn{4}{l}{\textit{AWP}} \\
In-degree & 0.884 & 0.728 & [0.718, 0.737] \\
In-value (sum inbound) & 0.669 & 0.495 & [0.481, 0.510] \\

\bottomrule
\end{tabular}
}
\end{table}


표~\ref{tab:alignment-transfer}에서, AWP와 입차수(들어오는 화살표 수)의 일치는 매칭 코호트에서 대리 지표별 켄달의 $\tau$ 가운데 가장 커요. 받은 양의 합과의 일치는 그보다 훨씬 낮은데, 그 까닭은 \ref{sec:same-window-alignment}절에 있어요. 엔도스랭크의 송금 계수는 중간 정도이고($0.341$과 $0.319$), 그 구간은 0과도, AWP의 구간과도 겹치지 않아요. 이 계수들은 주로, 허락 그래프 안의 지갑이 그 밖의 지갑보다 송금도 더 많이 받는다는 점을 보여 줘요(\ref{sec:rank-divergence}절). 평판을 들어오는 흐름과 같은 것으로 보는 서비스라면, 그 흐름은 입차수만 봐도 읽을 수 있어요. 그리고 기간 밖에서 AWP는 새 송금자를 입차수보다 못하게 줄 세워요(\ref{sec:holdout-results}절).

\dissertationSubsubheading[sub:allowance-family]{허락 가족}

% Real BigQuery data -- regenerate with:
%   2-Dissertation-Draft-Works/wallet-reputation-experiments/scripts/run_dissertation_eval.py --real --export-latex
% Generated by export_latex_results.py -- do not edit by hand unless re-export fails
\begin{table}[htbp]
\centering
\caption[Domain alignment with allowance proxies.]{Domain alignment with allowance proxies (spender-side in-approve degree and value). CI = 95\% percentile bootstrap over wallets (400 paired resamples, seed 42).}
\label{tab:alignment-allowance}
\fitwidth{%
\begin{tabular}{lccc}
\toprule
Proxy & Spearman $\rho$ & Kendall $\tau$ & 95\% CI of $\tau$ \\
\midrule
\multicolumn{4}{l}{\textit{EndorseRank}} \\
In-approve degree & 0.380 & 0.376 & [0.344, 0.403] \\
In-approve value (sum) & 0.380 & 0.373 & [0.341, 0.398] \\
\addlinespace
\multicolumn{4}{l}{\textit{AWP}} \\
In-approve degree & -0.029 & -0.024 & [-0.043, -0.002] \\
In-approve value (sum) & -0.029 & -0.023 & [-0.043, -0.002] \\

\bottomrule
\end{tabular}
}
\end{table}


표~\ref{tab:alignment-allowance}에서 엔도스랭크의 두 계수는 거의 같고, 동점 때문에 얻을 수 있는 가장 큰 값(천장)에 닿아 있어요(\ref{sec:same-window-alignment}절). AWP의 두 계수는 조금 음수이고($-0.024$와 $-0.023$), 구간은 0을 벗어나요. 곧 코호트 지갑들 가운데서, 송금 그래프에서 중심에 있는 정도(중심성)로는 스펜더로 허락받은 지갑을 골라내지 못해요. 이 비교는 허락을 받는 지갑 131개에 기대고 있고, 처음부터 어느 정도 생길 수밖에 없는 결과예요. 엔도스랭크가 바로 그 허락들로 계산되기 때문이에요.

\dissertationSubsubheading[sub:sybil-family]{시빌 보정 안정성 가족}

% Real BigQuery data -- regenerate with:
%   2-Dissertation-Draft-Works/wallet-reputation-experiments/scripts/run_dissertation_eval.py --real --export-latex
% Generated by export_latex_results.py -- do not edit by hand unless re-export fails
\begin{table}[htbp]
\centering
\caption[Domain alignment with Sybil-adjusted stability proxies from transfer activity.]{Domain alignment with Sybil-adjusted stability proxies from transfer activity. CI = 95\% percentile bootstrap over wallets (400 paired resamples, seed 42).}
\label{tab:alignment-sybil-stability}
\fitwidth{%
\begin{tabular}{lccc}
\toprule
Proxy & Spearman $\rho$ & Kendall $\tau$ & 95\% CI of $\tau$ \\
\midrule
\multicolumn{4}{l}{\textit{EndorseRank}} \\
Inbound counterparty ratio & 0.091 & 0.072 & [0.045, 0.095] \\
Transfer tenure (days) & 0.363 & 0.294 & [0.274, 0.316] \\
Active months & 0.377 & 0.327 & [0.307, 0.348] \\
\addlinespace
\multicolumn{4}{l}{\textit{AWP}} \\
Inbound counterparty ratio & -0.269 & -0.228 & [-0.250, -0.207] \\
Transfer tenure (days) & 0.685 & 0.507 & [0.493, 0.521] \\
Active months & 0.704 & 0.555 & [0.540, 0.569] \\

\bottomrule
\end{tabular}
}
\end{table}


표~\ref{tab:alignment-sybil-stability}를 보면, AWP는 활동 기간(첫 기록부터 마지막 기록까지 지난 날 수)에서 $0.507$ 대 $0.294$로, 활동한 달 수에서 $0.555$ 대 $0.327$로 엔도스랭크보다 앞서요. 기간 전체에 퍼진 활동에 상을 주는 송금 점수라면 그럴 것이라고 예상할 수 있어요. 가족 평균에 대한 짝지은 차이도 구간이 0을 벗어나요(표~\ref{tab:tau-diff}, 대비~iv). 들어오는 상대 비율(받은 송금 수에 견준 서로 다른 보낸 쪽 수)에서는 AWP가 뚜렷하게 음수이고($-0.228$), 엔도스랭크는 약하게 양수예요($0.072$). 이 가족은 각 점수가 시간에 걸친 송금 정보를 얼마나 담고 있는지 보여 줘요.

표~\ref{tab:alignment-summary}의 보관해 둔 거래 가족에서, AWP는 청산(손실이 너무 커져 거래가 강제로 끝나는 일)에서 $0.141$ 대 $0.093$으로, 거래 성공에서 $0.236$ 대 $0.096$으로 엔도스랭크보다 앞서요. 위험 반대(값이 클수록 손실이 적은 지표) 가족에서는 둘 다 음수로, 엔도스랭크는 $-0.055$, AWP는 $-0.127$이에요. 손실 없이 닫은 비율은 거래량이 늘어도 커지지 않는 비율인데, 여기서 두 계수는 $-0.031$과 $-0.006$이에요.

\dissertationSubheading[sec:er-awp-holdout]{기간 밖}

표~\ref{tab:endorserank-awp-holdout}는 두 홀드아웃 기간에서 두 점수 사이의 짝지은 차이를 모아 두었어요. 스펜더(spender, 허락을 받아 대신 쓰는 쪽) 코호트 전체에서는, 라벨(나중에 채점에 쓰는 정답 값)이 세는 화살표 종류로 만든 점수를 $a$로 삼아요. 두 번째 묶음의 부분 집합에서는 두 행 모두 엔도스랭크를 $a$로 삼아요.

% Spring holdout and post hoc June-August scores -- regenerate with:
%   2-Dissertation-Draft-Works/wallet-reputation-experiments/scripts/run_holdout_eval.py --cohort spenders
%   2-Dissertation-Draft-Works/wallet-reputation-experiments/scripts/run_fresh_posthoc.py
% Generated by export_latex_results.py -- do not edit by hand unless re-export fails
\begin{table}[htbp]
\centering
\caption[Paired bootstrap differences between EndorseRank and AWP out of window.]{Paired bootstrap differences between EndorseRank and AWP out of window. None was fixed in advance: the spring rows were computed after the spring labels were known, and EndorseRank was scored on the June--August cohorts after the registered evaluation (Table~\ref{tab:fresh-posthoc}). The second block is the subset of Table~\ref{tab:holdout-receiving}. Within each block $a$ and $b$ share the wallet resamples (400, seed 42).}
\label{tab:endorserank-awp-holdout}
\fitwidth{%
\begin{tabular}{llccccc}
\toprule
Contrast ($a$ minus $b$) & Label & $\tau_a$ & $\tau_b$ & $\Delta\tau$ & 95\% CI of $\Delta\tau$ & $P(\Delta\tau>0)$ \\
\midrule
\multicolumn{7}{l}{\emph{Spring holdout, spenders ($n=1{,}335$)}} \\
EndorseRank minus AWP & new approval pairs & 0.394 & 0.123 & +0.271 & [0.207, 0.332] & 1.000 \\
AWP minus EndorseRank & new transfer senders & 0.344 & 0.283 & +0.061 & [0.026, 0.098] & 1.000 \\
\midrule
\multicolumn{7}{l}{\emph{Spring holdout, spenders that had received a transfer before the freeze ($n=977$)}} \\
EndorseRank minus AWP & new approval pairs & 0.340 & 0.298 & +0.042 & [0.003, 0.085] & 0.978 \\
EndorseRank minus AWP & new transfer senders & 0.358 & 0.388 & -0.029 & [-0.070, 0.011] & 0.085 \\
\midrule
\multicolumn{7}{l}{\emph{June--August, spenders ($n=1{,}964$)}} \\
EndorseRank minus AWP & new approval pairs & 0.333 & 0.118 & +0.215 & [0.171, 0.270] & 1.000 \\
AWP minus EndorseRank & new transfer senders & 0.319 & 0.255 & +0.064 & [0.036, 0.088] & 1.000 \\
\midrule
\multicolumn{7}{l}{\emph{June--August, traders with at least three closes ($n=1{,}402$)}} \\
EndorseRank minus AWP & liquidation-free close rate & 0.081 & 0.045 & +0.035 & [-0.017, 0.087] & 0.925 \\
\bottomrule
\end{tabular}
}
\end{table}


봄 스펜더에서 엔도스랭크는 새 허락 쌍(새로 생긴 주인–스펜더 허락)을 $\tau=0.394$(구간 $[0.360,0.431]$)로 줄 세우고, AWP는 $0.123$이에요. 짝지은 차이는 $+0.271$($[0.207,0.332]$)이에요. 새 송금자(새로 송금을 보내온 주소)에서는 순서가 뒤집혀요. AWP는 $0.344$이고 엔도스랭크는 $0.283$이어서, 차이는 $+0.061$($[0.026,0.098]$)이에요. 모든 라벨은 점수를 매긴 모든 이벤트보다 나중에 일어나요. 그래서 이 코호트에서는 화살표 종류마다, 다른 쪽에는 없는 자기 구성개념에 대한 정보를 담고 있어요. 첫 번째 차이의 많은 부분은 다루는 범위로 설명돼요. AWP는 동결일 전에 송금 화살표가 없는 스펜더 318개에 0점을 주는데, 이 스펜더들은 나머지보다 새 허락 쌍을 더 자주 얻었어요. 송금을 받은 적이 있는 스펜더 977개에서는, 새 허락 쌍에서 엔도스랭크가 앞서는 폭이 $+0.042$($[0.003,0.085]$)로 줄어요. 그리고 새 송금자에서는 두 점수가 비슷해요($-0.029$, $[-0.070,0.011]$). 송금 정보를 전혀 쓰지 않는 받은 허락 수(허락 입차수)도 새 송금자에서 AWP보다 앞서요($0.404$ 대 $0.344$, 짝지은 검사는 하지 않음).

6월부터 8월까지의 기간에서도 같은 모습이 되풀이돼요. 이 기간에는 등록해 둔 평가를 마친 뒤에 엔도스랭크 점수를 매겼어요. 엔도스랭크는 새 허락 쌍을 $0.333$으로, AWP는 $0.118$로 줄 세워요($+0.215$, $[0.171,0.270]$). AWP는 새 송금자를 $0.319$로, 엔도스랭크는 $0.255$로 줄 세워요($+0.064$, $[0.036,0.088]$). 트레이더(거래하는 사람)의 청산 없이 닫은 비율에서는 엔도스랭크가 $0.081$, AWP가 $0.045$예요. 이 차이는 0과 구별할 수 없어요($+0.035$, $[-0.017,0.087]$).

탐색적(미리 정하지 않고 살펴본) 봄 실행의 매칭 트레이더(표~\ref{tab:holdout-traders})에서, AWP는 이익을 낸 닫기 횟수를 $0.210$(구간 $[0.181,0.238]$)으로 줄 세우고 엔도스랭크는 $0.072$($[0.041,0.111]$)예요. 실현한 이익은 AWP가 $0.191$($[0.159,0.219]$), 엔도스랭크가 $0.048$($[0.017,0.082]$)이에요. 이익을 낸 닫기의 비율에서는 어느 점수도 아무것도 보여 주지 못해요. AWP는 $0.016$($[-0.015,0.044]$)이고 엔도스랭크는 $-0.016$($[-0.053,0.024]$)이에요. AWP가 앞서는 것은 횟수와 합계인데, 이것들은 지갑이 거래를 많이 할수록 커져요. 그리고 같은 두 라벨에서 송금 입차수도 AWP와 거의 비슷하게 잘해요.
